\chapterdivider{8}{Two-parameter fold and Hopf curves}
  {How do local bifurcations move when two physical parameters vary?}
  {Featured application: trace a fold curve through a Bogdanov--Takens point}

\begin{frame}{From one organizing point to a curve}
\releasedbadge
\small
\begin{columns}[T,onlytextwidth]
\column{.47\textwidth}
\begin{block}{Direct point refinement}
Supply a nearby state and parameter pair; solve the complete codimension-two
extended system for one isolated CP, BT, GH, ZH, or HH point.
\end{block}
\column{.49\textwidth}
\begin{block}{Two-parameter curve continuation}
Supply one checked fold or Hopf seed; continue its codimension-one defining
system while a second physical parameter varies.
\end{block}
\end{columns}
\vfill
\begin{alertblock}{Different jobs, composable workflow}
A direct solver can refine a candidate organizer. A fold/Hopf curve then shows
how the surrounding local event moves through the parameter plane.
\end{alertblock}
\end{frame}

\begin{frame}{One physical parameter becomes part of the state}
\small
For an equilibrium model $f(u,p_0,p_1)=0$, select one component
$q=p_{\rm free}$ as the continuation coordinate. The other physical
parameter is solved inside an extended state.
\vspace{3mm}
\begin{columns}[T,onlytextwidth]
\column{.48\textwidth}
\begin{block}{Fold curve}
\[
X_{\rm fold}=(u,p_{\rm fixed},v),
\qquad q=p_{\rm free}.
\]
The residual contains $f=0$, $f_u v=0$, and a normalization for the fold
null vector $v$.
\end{block}
\column{.48\textwidth}
\begin{block}{Hopf curve}
\[
X_{\rm H}=(u,p_{\rm fixed},q_1,q_2,\omega).
\]
The residual contains $f=0$, the real/imaginary eigenvector equations, and
phase/normalization constraints.
\end{block}
\end{columns}
\vfill
\begin{alertblock}{The reuse is exact}
To the continuation engine this is again a root $F(X,q)=0$ in one scalar
parameter. The same fixed-buffer pseudo-arclength machinery can trace it.
\end{alertblock}
\end{frame}

\begin{frame}{Choose which parameter moves}
\centering
\begin{tikzpicture}[node distance=6mm and 8mm,>=Latex,
 box/.style={rounded corners,draw=JaxBlue,thick,fill=SoftBlue,
             minimum height=9mm,minimum width=42mm,align=center},
 result/.style={rounded corners,draw=JaxGreen,thick,fill=JaxGreen!10,
                minimum height=9mm,minimum width=48mm,align=center}]
  \node[box] (physical) {$p=(p_0,p_1)$\\two physical parameters};
  \node[box,right=of physical] (choice) {\texttt{free=1}\\continue $q=p_1$};
  \node[result,right=of choice] (packed) {$p_0$ solved inside $X$\\$p_1$ stored as branch parameter};
  \draw[->,very thick,JaxTeal] (physical)--(choice);
  \draw[->,very thick,JaxTeal] (choice)--(packed);
\end{tikzpicture}
\vspace{5mm}
\begin{block}{Two values must agree}
If the factory receives \texttt{p\_guess = [p0, p1]}, then
\texttt{continuation(..., p\_span=(p1, end))} must start at
\texttt{p\_guess[free]}. A mismatch is rejected instead of silently moving
the initial point.
\end{block}
\begin{alertblock}{Shape and index contract}
\texttt{p\_guess.shape == (2,)} and \texttt{free} is either 0 or 1. These
factories describe codimension-two parameter geometry, not an arbitrary
number of free parameters.
\end{alertblock}
\end{frame}

\begin{frame}[fragile]{Build and continue a fold curve}
\small
\begin{lstlisting}
prob = jc.fold_curve_problem(
    rhs,
    u_guess,
    jnp.array([p0_guess, p1_guess]),
    free=1,
    args=args,
)

result = jc.continuation(
    prob,
    p_span=(p1_guess, p1_end),
    settings=jc.ContinuationPar(
        compute_stability=False, newton_tol=1e-5
    ),
    events=[jc.BogdanovTakens(
        raw_f=rhs, free=1, curve="fold", args=args
    )],
)
\end{lstlisting}
\begin{block}{Two Newton roles}
The factory first refines the guess with the fold extended system. During the
sweep, the ordinary continuation corrector advances the already packed curve
state. Both convergence checks matter.
\end{block}
\end{frame}

\begin{frame}[fragile]{The Hopf-curve factory has the same outer contract}
\small
\begin{lstlisting}
prob = jc.hopf_curve_problem(
    rhs, u_guess, p_guess, free=0, args=args
)

result = jc.continuation(
    prob,
    p_span=(p_guess[0], p0_end),
    settings=jc.ContinuationPar(
        compute_stability=False, newton_tol=1e-5
    ),
    events=[
        jc.GeneralizedHopf(raw_f=rhs, free=0, args=args),
        jc.ZeroHopf(raw_f=rhs, free=0, curve="hopf",
                    args=args),
    ],
)
\end{lstlisting}
\begin{alertblock}{Why \texttt{compute\_stability=False}?}
The stored state is an extended fold/Hopf system, not the original physical
state. Curve events evaluate the original Jacobian through \texttt{raw\_f};
generic stability of the extended state would answer the wrong question.
\end{alertblock}
\end{frame}

\begin{frame}{Which codimension-two events live on which curve?}
\small
\centering
\renewcommand{\arraystretch}{1.25}
\begin{tabular}{p{.17\textwidth}p{.28\textwidth}p{.38\textwidth}}
\toprule
Curve & Available event classes & Additional condition detected \\
\midrule
Fold & CP, BT, ZH & fold coefficient $a=0$; a second zero mode; or an
additional imaginary pair \\
Hopf & BT, GH, ZH, HH & $\omega\to0$; $l_1=0$; an additional zero mode; or a
second distinct imaginary pair \\
\bottomrule
\end{tabular}
\vspace{4mm}
\begin{block}{Events decode physical coordinates}
Every curve event carries the original \texttt{raw\_f}, the \texttt{free}
index, optional model \texttt{args}, and an explicit \texttt{curve} kind.
Refined hits store the physical parameter vector in \texttt{hit.info["p"]}.
\end{block}
\begin{alertblock}{Do not infer the curve type from a branch point}
The generic event protocol sees packed arrays, not problem metadata. Explicit
curve identity prevents applying an equilibrium test to the wrong residual.
\end{alertblock}
\end{frame}

\begin{frame}{Discard the condition that defines the curve}
\small
\begin{columns}[T,onlytextwidth]
\column{.47\textwidth}
\begin{block}{Along a fold curve}
One eigenvalue is pinned at zero everywhere. Detecting BT or ZH begins by
removing the entry nearest zero, then asking whether an \emph{additional}
real-zero or imaginary-pair condition appears.
\end{block}
\column{.49\textwidth}
\begin{block}{Along a Hopf curve}
The pair nearest $\pm i\omega$ is pinned to the imaginary axis everywhere.
Detecting ZH or HH removes that pair before searching the remaining spectrum.
Near BT, both collapsing entries nearest zero must be removed.
\end{block}
\end{columns}
\vspace{4mm}
\begin{alertblock}{Otherwise every point looks special}
The defining zero mode or imaginary pair contains no new codimension-two
information. Counting it again creates false candidates rather than events.
\end{alertblock}
\end{frame}

\begin{frame}{Example 15: an exact Bogdanov--Takens geometry}
\small
For shifted variables $x=X-5$, $y=Y-2$, define
\[
x'=y,\qquad
y'=\beta_1+\beta_2x+x^2+xy,
\]
with $\beta_1=p_0-3$ and $\beta_2=p_1+1$.
\begin{columns}[T,onlytextwidth]
\column{.47\textwidth}
\begin{block}{Exact fold curve}
At an equilibrium $y=0$. The fold conditions give
\[
p_0=3+\frac{(p_1+1)^2}{4}.
\]
The double-zero point is exactly $(p_0,p_1)=(3,-1)$.
\end{block}
\column{.49\textwidth}
\begin{block}{Numerical experiment}
Start from $(3.25,-2)$, trace with \texttt{free=1}, and attach a
\texttt{BogdanovTakens(..., curve="fold")} event. The refined hit can be
compared directly with the exact organizer.
\end{block}
\end{columns}
\end{frame}

\begin{frame}{Example 15: read the parameter-plane diagram}
\centering
\small
\begin{tikzpicture}[x=2.2cm,y=1.1cm,>=Latex]
  \draw[->] (-2.25,0.0)--(.45,0.0) node[right] {$p_1$ (free)};
  \draw[->] (-2.1,-.2)--(-2.1,2.45) node[above] {$p_0$ (solved)};
  \draw[JaxGreen,very thick,domain=-2:0.25,samples=80]
    plot (\x,{3 + (\x+1)*(\x+1)/4 - 2.75});
  \filldraw[JaxOrange] (-1,.25) circle (2.7pt)
    node[above right] {BT: $(p_1,p_0)=(-1,3)$};
  \node[JaxGreen] at (-.05,1.15) {fold curve};
\end{tikzpicture}
\vspace{1mm}
\begin{block}{Plot the reconstructed physical coordinates}
\centering
\texttt{jc.plot\_two\_parameter\_diagram([(result, "fold")], free=1)}
\end{block}
\begin{alertblock}{Axes follow \texttt{free}}
The horizontal axis is $p_{\rm free}$ and the vertical axis is the parameter
decoded from the packed state. Pass the same \texttt{free} used by the factory.
\end{alertblock}
\end{frame}

\begin{frame}[fragile]{Batch curves without rebuilding the factory}
\small
\begin{lstlisting}
base = jc.fold_curve_problem(
    rhs_shiftable, u_guess, p_guess,
    free=1, args=0.0,
)

def endpoint(shift):
    out = jc.continuation(
        base.at(args=shift), p_span=(-2.0, 0.0),
        settings=settings,
    )
    last = jnp.sum(out.branch.valid) - 1
    return jnp.take(out.branch.states, last, axis=0)[2]

endpoints = jax.vmap(endpoint)(shifts)
\end{lstlisting}
\begin{alertblock}{Two transformation rules}
Build the seed-refining factory outside \texttt{vmap}; carry swept values
through \texttt{args}. In traced output, use \texttt{Branch.valid} rather than
reading padded buffer slots as real curve points.
\end{alertblock}
\end{frame}

\begin{frame}[fragile]{Differentiate a codimension-two location}
\small
\begin{lstlisting}
def bt_p0(shift):
    p_star = jc.bogdanov_takens_parameters(
        rhs_shiftable, u_guess, p_guess,
        args=shift,
    )
    return p_star[0]

sensitivity = jax.grad(bt_p0)(0.0)
\end{lstlisting}
\begin{columns}[T,onlytextwidth]
\column{.47\textwidth}
\begin{block}{What is differentiated?}
The converged extended-system equation is differentiated by the implicit
function theorem. Solver iterations are not replayed in reverse.
\end{block}
\column{.49\textwidth}
\begin{alertblock}{Check before composing}
First call and verify the corresponding \texttt{bogdanov\_takens\_point}
result. The parameter-only wrapper omits the convergence flag to remain
grad-ready.
\end{alertblock}
\end{columns}
\end{frame}

\begin{frame}{Practical limits and restart strategy}
\small
\begin{columns}[T,onlytextwidth]
\column{.47\textwidth}
\begin{block}{Supported in v0.4.0}
\begin{itemize}
  \item equilibrium fold and Hopf curves;
  \item CP/BT/GH/ZH/HH events where applicable;
  \item parameter-plane plots;
  \item fixed-shape JIT and batching;
  \item implicit sensitivities of direct organizers.
\end{itemize}
\end{block}
\column{.49\textwidth}
\begin{alertblock}{Boundaries}
\begin{itemize}
  \item no PD/LPC/NS periodic-orbit curves;
  \item no branch switching at an organizer;
  \item no automatic discovery of disconnected curves;
  \item no adaptive Hopf eigenvector re-anchoring inside a traced buffer.
\end{itemize}
\end{alertblock}
\end{columns}
\vspace{2mm}
If the Hopf eigenvector rotates toward orthogonality with its initial anchor,
restart from a later checked point with a fresh local anchor.
\end{frame}

\begin{frame}{What to remember}
\begin{enumerate}
  \item A two-parameter fold/Hopf curve is an extended root system continued
        in one selected physical parameter.
  \item The factory refines the seed; \texttt{p\_span[0]} must match the
        selected component of the supplied parameter pair.
  \item Curve events remove the eigenvalue condition already pinned by the
        curve before testing for an additional degeneracy.
  \item Build factories outside JAX transforms, pass sweeps through
        \texttt{args}, and respect traced validity masks.
\end{enumerate}
\vfill
\begin{alertblock}{Scientific interpretation}
The curve organizes local equilibrium bifurcations in the chosen parameter
plane. It does not by itself construct the branches or cycles that emanate
from those events.
\end{alertblock}
\end{frame}
